% Technical note for collaborators: Risk-Indexed Threshold (RIT) versus
% Grid-Searched Threshold (GST) for cached dispatch (the s-only pivot). Companion to dispatch_note.tex and redundancy_note.tex.
% ICLR 2027 template, no title block. Compile with pdflatex.
\documentclass{article}
\usepackage{iclr2027_conference,times}
\iclrfinalcopy
\usepackage{amsmath,amssymb}
\usepackage{enumitem}

\begin{document}

\paragraph{Problem.}
Each decision step of the cached executor observes a single scalar, the fused
retrieval score $s_t$, and must choose one of $K{+}1$ actions: full inference
(action $0$) or one of $K$ reuse tiers $a \in \{1, \dots, K\}$ of increasing
skip depth $\tau_1 < \dots < \tau_K$ (in the current system $K = 2$: warm
start at $\mathtt{start\_t} = 0.3$ and full replay). Unit compute costs are
fixed and known,
\begin{equation}
c_a \;=\; \mathrm{s}_1 \;+\; \mathbf{1}[a < K]\,\bigl( \mathrm{s}_2 +
\mathtt{start\_t}(a)\cdot \mathrm{s}_3 \bigr),
\qquad
\mathtt{start\_t}(0) = 1, \quad \mathtt{start\_t}(K) = 0,
\label{eq:cost}
\end{equation}
so that $c_K < \dots < c_1 < c_0 = \mathrm{s}_1 + \mathrm{s}_2 + \mathrm{s}_3$,
with $\mathrm{s}_1, \mathrm{s}_2, \mathrm{s}_3$ the measured stage latencies
(vision encoding, language-model prefix, full flow integration). Every
decision pays $\mathrm{s}_1$, which produces the features the retrieval score
is computed from; a warm tier resumes the flow at $\mathtt{start\_t}(a)$ and
pays that fraction of $\mathrm{s}_3$; the deepest tier replays the cached
chunk and skips both the prefix and the integration.
The online rule class we consider is the incumbent one: \emph{monotone cut
vectors}. A rule is a vector $\theta = (\theta_1, \dots, \theta_K)$,
$\theta_1 \le \dots \le \theta_K$ interpreted through ``execute the deepest
tier whose cut is cleared'', $a(s) = \max(\{0\} \cup \{a : s \ge
\theta_a\})$. Two ways of choosing $\theta$ are compared throughout. The incumbent, which
we call the \emph{Grid-Searched Threshold} (GST), sets the cut vector at
development-score percentiles and finds operating points by rolling out a
grid of them; the alternative, the \emph{Risk-Indexed Threshold} (RIT),
computes the cut vector from calibrated risk curves at one tolerance. Both
are members of this class at run time and differ only in \emph{how $\theta$
is chosen offline}. This note makes that difference the object of study.

\paragraph{Two indexings of one family.}
GST is \emph{rate indexing}: it sets each $\theta_a$ at a percentile of
the development score distribution: $\theta_a = \hat{F}_s^{-1}(1 - f_a)$
admits a target fraction $f_a$ of decisions into tier $a$ or deeper. Rates
are predictable by construction; \emph{risk} is not. Which vector
$(f_1, \dots, f_K)$ lands on the cost--success frontier is unknown a priori,
so GST finds operating points by rolling out a $K$-dimensional grid. RIT
is \emph{risk indexing}. Let $Y_a \ge 0$
be the offline-checkable deviation surrogate of the companion note
(Eq.~2 there) for tier $a$, and define per tier the conditional upper
quantile
\begin{equation}
q_a(s) \;=\; Q_{1-\alpha}\bigl( Y_a \mid s_t = s \bigr),
\qquad a = 1, \dots, K,
\label{eq:qa}
\end{equation}
estimated on calibration rollouts by order-restricted (isotonic) quantile
regression under the monotonicity assumption (A1$'$): $q_a$ nonincreasing in
$s$, nondecreasing in the skip depth. Given one tolerance $\delta$, the cuts
are \emph{computed}, not searched:
\begin{equation}
s^{*}_a(\delta) \;=\; \min\{\, s : \hat{q}_a(s) \le \delta \,\},
\qquad
\theta_a = s^{*}_a(\delta) .
\label{eq:cuts}
\end{equation}
The rule executes the deepest admissible tier and falls back to full
inference when no tier is admissible (fail-closed, as in the companion
note). Online, Eq.~\ref{eq:cuts} is indistinguishable in form and cost from
GST: the same comparisons against stored constants.

\paragraph{Equivalence at $K = 1$ (a concession).}
With a single tier the two indexings coincide as families. Under (A1$'$) the
map $\delta \mapsto s^{*}_1(\delta)$ is a monotone reparametrization of the
cut, so
$\{\text{RIT rules}\}_{\delta} =
\{\text{GST rules}\}_{\theta} =
\{\text{monotone rules on } s\}$,
and under the monotone-likelihood-ratio premise already stated in the
companion note the common family traces the ROC frontier of $s$
(Neyman--Pearson). At $K=1$, therefore, RIT buys semantics but no
new operating points, and percentile inversion is equally automatic. Any
claimed contribution must---and here does---begin at $K \ge 2$.

\paragraph{Coupling at $K \ge 2$: the solution path.}
Consider the budgeted design problem over cut vectors,
\begin{equation}
\max_{\theta}\ \mathrm{SR}(\theta)
\qquad \text{s.t.} \qquad
\mathbb{E}\bigl[ c_{a(s;\theta)} \bigr] \le B .
\label{eq:budget}
\end{equation}
Introduce the surrogate objective in which closed-loop degradation is
exchanged for accumulated per-decision surrogate risk,
\begin{equation}
\text{(A2)} \qquad
\mathrm{SR}(\theta) \;\approx\; \Phi\Bigl(
\mathbb{E}\bigl[\, r_{a(s;\theta)}(s) \,\bigr] \Bigr),
\qquad
r_a(s) = \text{a nondecreasing functional of } q_a(s),\ r_0 \equiv 0,
\label{eq:A2}
\end{equation}
with $\Phi$ decreasing. (A2) is an approximation, not a theorem: it ignores
sequencing and closed-loop intervention (the composition failure of the
companion note) and is adopted because its consequences are checkable. Under
(A1$'$), (A2), and $c_K < \dots < c_0$, the Lagrangian of
Eq.~\ref{eq:budget} separates across decisions and across adjacent tier
boundaries: at an optimum, each boundary $\theta_a$ equalizes marginal
surrogate risk against marginal compute saving with the \emph{same}
multiplier $\lambda$,
\begin{equation}
\frac{\partial\, \mathbb{E}[r]}{\partial \theta_a}
\;=\;
\lambda\, \frac{\partial\, \mathbb{E}[c]}{\partial \theta_a}
\qquad (a = 1, \dots, K),
\label{eq:kkt}
\end{equation}
so the set of optimal cut vectors is a \emph{one-parameter solution path}
indexed by $\lambda$. When the marginal risk at a boundary is governed by
the calibrated quantile at the cut---the case in which admission at
$\theta_a$ adds decisions of surrogate risk $\approx \hat q_a(\theta_a)$---
the path is traced by a common level: $\hat q_a(\theta_a) = \delta(\lambda)$
for all $a$, which is exactly Eq.~\ref{eq:cuts}. The RIT ladder
$\{\,\theta(\delta) : \delta \ge 0\,\}$ is thus the plug-in estimate of the
solution path of the $K$-dimensional design problem, and GST's
$K$-dimensional grid search rediscovers that path pointwise. Two consequences are testable:
(i) RIT ladder points should track GST's empirical Pareto hull at
matched evaluation budget; (ii) at a fixed number $n$ of a-priori arms, the
vertical regret of GST's $K$-dimensional grid against its own hull grows with
$K$, while RIT remains one-dimensional. Neither consequence asserts
that RIT \emph{dominates} GST---it cannot, being a subset;
the claim is that the path attains the family's frontier without the search.

\paragraph{Per-task and cross-library transfer.}
GST conditions on nothing: a percentile vector tuned on one task
mixture or one library encodes that mixture's score distribution, and the
GST's own documentation records that carrying a threshold across
libraries is a known failure mode. RIT recomputes
Eq.~\ref{eq:cuts} from the calibrated curves: per-task (Mondrian) cuts
$s^{*}_{a,t}(\delta)$ follow from per-task quantile curves at the same
$\delta$ when sample sizes permit, and a library change re-runs calibration
rather than a $K$-dimensional (per task, $K \times T$-dimensional) grid.
The empirical finding that task-conditional selection improves every family
by a comparable margin is consistent with this reading: conditioning is a
property of the information set, not of any single rule family; what RIT
changes is the \emph{cost of exploiting it}.

\paragraph{What is not claimed.}
No statement here is a closed-loop guarantee: deployment intervenes on the
visitation distribution, and Eq.~\ref{eq:A2} is a surrogate adopted for its
checkable consequences (the companion note's scope paragraph applies
verbatim). The path property is relative to rules measurable in $s$ alone;
richer information sets---the disagreement statistic $v$, task identity used
online, stateful gates---can move the frontier, and whether they do is an
empirical question that our measurements answer in the negative for $v$
(no measurable increment at matched density) and in the positive, equally
for all families, for task identity. Finally, the estimated path can miss
the frontier where calibration data are sparse; matched-density evaluation,
not construction, is what supports the tracking claim.

\paragraph{Experimental consequences.}
Four readouts follow. (i) \emph{Matched-density frontier comparison}: the
RIT ladder against the dense GST grid on the same evaluation set,
library, and cost accounting; the tracking claim predicts hull agreement
within estimation noise. (ii) \emph{Tuning budget}: frontier quality (hull
area or vertical regret) as a function of the number of a-priori arms for
each method; the prediction is a widening gap as arms are removed, and a
qualitative jump when $K$ grows. (iii) \emph{Risk calibration}: realized
per-decision surrogate risk of admitted steps against the declared $\delta$
---the plot only RIT can draw in advance. (iv) \emph{System
composition}: the full deployed stack adds the production hysteresis gate in
front of either method; the ceiling-equality reading predicts that gate
plus RIT matches gate plus GST at matched freedom, measured
on the same chain. A three-tier instance ($K = 3$, second warm depth
unpinned from $\mathtt{start\_t}=0.3$ after its unit cost is frozen) turns
consequence (ii) into the discriminating experiment: budget-matched
$3$-dimensional GST grid against the unchanged one-dimensional RIT ladder.

\end{document}
